% =============================================================
% 02_management.tex  第2部 マネジメント系
% =============================================================
\chapter{マネジメント系：どう作り、どう回すか}

\section{システム開発工程}

\pitfall{工程名がカタカナだらけで頭に入ってこない}{
開発工程の用語は、似た意味で呼び方だけ違うものが多く存在します。
まず\textbf{5 つのコア工程}を固定し、それ以外はバリエーションとして覚えるのが近道です。
}

\figurebox{システム開発工程の流れ}{fig:sdlc}{%
  % 図: システム開発工程の流れ図
% 用途: % 図: システム開発工程の流れ図
% 用途: % 図: システム開発工程の流れ図
% 用途: \input{assets/diagrams/sdlc_flow}
\begin{tikzpicture}[node distance=4mm and 5mm]

  \node[boxnode, minimum width=19mm] (req)  {要件定義};
  \node[boxnode, minimum width=19mm, right=of req]  (design) {設計};
  \node[boxnode, minimum width=19mm, right=of design] (impl)  {実装};
  \node[boxnode, minimum width=19mm, right=of impl]   (test)  {テスト};
  \node[boxnode, minimum width=19mm, right=of test]   (op)    {運用・保守};

  \draw[arr] (req)    -- (design);
  \draw[arr] (design) -- (impl);
  \draw[arr] (impl)   -- (test);
  \draw[arr] (test)   -- (op);

  \node[subnode, below=5mm of req,    minimum width=19mm] {何を作る};
  \node[subnode, below=5mm of design, minimum width=19mm] {どう作る};
  \node[subnode, below=5mm of impl,   minimum width=19mm] {作る};
  \node[subnode, below=5mm of test,   minimum width=19mm] {動くか};
  \node[subnode, below=5mm of op,     minimum width=19mm] {使い続ける};

  \node[subnode, below=14mm of impl, minimum width=108mm, text width=100mm, align=center]
    {ウォーターフォール：一方向に進む／アジャイル：この流れを小さく何周もする};

\end{tikzpicture}

\begin{tikzpicture}[node distance=4mm and 5mm]

  \node[boxnode, minimum width=19mm] (req)  {要件定義};
  \node[boxnode, minimum width=19mm, right=of req]  (design) {設計};
  \node[boxnode, minimum width=19mm, right=of design] (impl)  {実装};
  \node[boxnode, minimum width=19mm, right=of impl]   (test)  {テスト};
  \node[boxnode, minimum width=19mm, right=of test]   (op)    {運用・保守};

  \draw[arr] (req)    -- (design);
  \draw[arr] (design) -- (impl);
  \draw[arr] (impl)   -- (test);
  \draw[arr] (test)   -- (op);

  \node[subnode, below=5mm of req,    minimum width=19mm] {何を作る};
  \node[subnode, below=5mm of design, minimum width=19mm] {どう作る};
  \node[subnode, below=5mm of impl,   minimum width=19mm] {作る};
  \node[subnode, below=5mm of test,   minimum width=19mm] {動くか};
  \node[subnode, below=5mm of op,     minimum width=19mm] {使い続ける};

  \node[subnode, below=14mm of impl, minimum width=108mm, text width=100mm, align=center]
    {ウォーターフォール：一方向に進む／アジャイル：この流れを小さく何周もする};

\end{tikzpicture}

\begin{tikzpicture}[node distance=4mm and 5mm]

  \node[boxnode, minimum width=19mm] (req)  {要件定義};
  \node[boxnode, minimum width=19mm, right=of req]  (design) {設計};
  \node[boxnode, minimum width=19mm, right=of design] (impl)  {実装};
  \node[boxnode, minimum width=19mm, right=of impl]   (test)  {テスト};
  \node[boxnode, minimum width=19mm, right=of test]   (op)    {運用・保守};

  \draw[arr] (req)    -- (design);
  \draw[arr] (design) -- (impl);
  \draw[arr] (impl)   -- (test);
  \draw[arr] (test)   -- (op);

  \node[subnode, below=5mm of req,    minimum width=19mm] {何を作る};
  \node[subnode, below=5mm of design, minimum width=19mm] {どう作る};
  \node[subnode, below=5mm of impl,   minimum width=19mm] {作る};
  \node[subnode, below=5mm of test,   minimum width=19mm] {動くか};
  \node[subnode, below=5mm of op,     minimum width=19mm] {使い続ける};

  \node[subnode, below=14mm of impl, minimum width=108mm, text width=100mm, align=center]
    {ウォーターフォール：一方向に進む／アジャイル：この流れを小さく何周もする};

\end{tikzpicture}
%
}

\definitionbx{開発工程の 5 ステップ}{
\termtag{要件定義} → \termtag{設計} → \termtag{実装} → \termtag{テスト} → \termtag{運用・保守}。
上流にいくほど\textbf{抽象度が高い}。上流のミスは下流で取り返せない、というのが原則。
}

\comparisonbx{ウォーターフォールとアジャイル}{%
\comptable{観点 & ウォーターフォール & アジャイル}{
進め方 & 一方向に一巡 & 小さく何周も \\
強み & 計画しやすい・契約しやすい & 変更に強い・学習が早い \\
弱み & 変更が高価 & 全体の見通しが出にくい \\
向く案件 & 要件が固い & 要件が動く
}
}

\pitfall{アジャイルは「計画しないこと」ではない}{
アジャイル＝無計画という誤解は危険です。
アジャイルは「\termtag{小さい計画を何度も立てる}」方法であって、計画を捨てているわけではありません。
イテレーション（\termtag{スプリント}）ごとに計画・実装・レビュー・振り返りを回します。
}

\section{プロジェクトマネジメント}

\definitionbx{プロジェクト}{
\textbf{期限}と\textbf{目的}が決まった、\textbf{一度きりの}活動。
日常の定常業務（オペレーション）とは区別される。
}

\comparisonbx{QCD（プロジェクトの三制約）}{%
\comptable{制約 & 中身 & 現実の葛藤}{
Q: 品質 (Quality) & 要求を満たす度合い & 落とすと手戻り \\
C: コスト (Cost) & 人件費・費用 & 削ると品質に跳ねる \\
D: 納期 (Delivery) & スケジュール & 縮めると人/品質に跳ねる
}
}

\advice{%
QCD は\textbf{「三角形のどれか一つを動かすと、他の二つが歪む」}と覚えます。
試験では「品質を守りつつコストを下げるには？」のような、三者のバランス感覚を問う問題が出ます。
}

\section{サービスマネジメント}

\pitfall{「開発」が終わればゴールだと思っている}{
実務では、リリース後の方がはるかに長く、コストもかかります。
ITパスポートではここを\termtag{サービスマネジメント}として問います。
主役は ITIL と SLA、それから問い合わせの窓口である\termtag{サービスデスク}です。
}

\comparisonbx{SLA と SLM}{%
\comptable{用語 & 意味 & 役割}{
SLA & サービス品質の\textbf{約束(契約)} & 何を守るかを決める \\
SLM & SLA を\textbf{管理し改善する活動} & 守れているかを運用で確認
}
}

\checkquestion{ウォーターフォールモデルの特徴として最も適切なものはどれか。}
  {要件変更に柔軟に対応しやすい。}
  {開発を小さな反復単位で進める。}
  {工程を上流から下流へ一方向に進める。}
  {設計と実装を同時並行で行う。}
  {C}
  {A と B はアジャイルの特徴、D はアジャイル寄りの開発現場の実態に近い記述です。ウォーターフォールの\textbf{本質は「逆流しないこと」}です。}

\onelinesummary{%
マネジメント系は\textbf{「どう作り、どう回すか」}。
開発工程・プロジェクト・サービス運用の 3 本柱を、\termtag{QCD という共通の三角形}で見ると一気に景色がつながる。
}
